% Part: applied-modal-logics
% Chapter: epistemic-logic
% Section: truth-at-w

\documentclass[../../../include/open-logic-section]{subfiles}

\begin{document}

\olfileid{aml}{el}{trw}

\olsection{কোনো জগতে সত্যতা}

স্বাভাবিক মোডাল যুক্তিবিদ্যার মতোই, প্রতিটি জ্ঞানতাত্ত্বিক মডেল নির্ধারণ করে তার কোন জগতে কোন !!{formula}s সত্য বলে গণ্য হবে। সম্পর্কভিত্তিক অর্থতত্ত্বের মৌলিক ধারণাটির জন্য একই সংকেত ব্যবহার করি: ``মডেল~$\mModel{M}$ জগৎ~$w$-তে !!{formula}~$!A$-কে সত্য করে''। সম্পর্কটি আরোহভাবে সংজ্ঞায়িত; মোডাল অপারেটরবিহীন সব অপারেটরের ক্ষেত্রে সংজ্ঞা স্বাভাবিক মোডাল যুক্তিবিদ্যার মতোই।

\begin{defn}\ollabel{defn:mmodels}
  $\mModel M$-এ $w$-তে !!a{formula}~$!A$-এর \emph{সত্যতা}, সংকেতে
  $\mSat{M}{!A}[w]$, নিম্নরূপে আরোহভাবে সংজ্ঞায়িত:
  \begin{enumerate}
  \tagitem{prvFalse}{\indcase{!A}{\lfalse}{$\mSat{M}{\lfalse}[w]$ কখনোই নয়।}}{}
  \tagitem{prvTrue}{\indcase{!A}{\ltrue}{$\mSat{M}{\ltrue}[w]$ সর্বদাই।}}{}
  \item $\mSat{M}{p}[w]$ তখনই এবং কেবল তখনই, যখন $w \in V(p)$।
  \tagitem{prvNot}{\indcase{!A}{\lnot !B}{$\mSat{M}{\indfrm}[w]$ তখনই এবং কেবল তখনই, যখন
    $\mSat/{M}{!B}[w]$।}}{}
  \tagitem{prvAnd}{\indcase{!A}{(!B \land !C)}{$\mSat{M}{\indfrm}[w]$ তখনই এবং কেবল তখনই, যখন
    $\mSat{M}{!B}[w]$ এবং $\mSat{M}{!C}[w]$।}}{}
  \tagitem{prvOr}{\indcase{!A}{(!B \lor !C)}{$\mSat{M}{\indfrm}[w]$ তখনই এবং কেবল তখনই, যখন
    $\mSat{M}{!B}[w]$ অথবা $\mSat{M}{!C}[w]$} (উভয়ই হতে পারে)।}{}
  \tagitem{prvIf}{\indcase{!A}{(!B \lif !C)}{$\mSat{M}{\indfrm}[w]$ তখনই এবং কেবল তখনই, যখন
    $\mSat/{M}{!B}[w]$ অথবা $\mSat{M}{!C}[w]$।}}{}
  \tagitem{prvIff}{\indcase{!A}{(!B \liff !C)}{$\mSat{M}{\indfrm}[w]$ তখনই এবং কেবল তখনই, যখন
    $\mSat{M}{!B}[w]$ ও $\mSat{M}{!C}[w]$ উভয়ই সত্য,
    অথবা $\mSat{M}{!B}[w]$ ও $\mSat{M}{!C}[w]$ কোনোটিই সত্য নয়।}}{}
  \item\ollabel{defn:sub:mmodels-box}
    \indcase{!A}{\Knows_a !B}{$\mSat{M}{\indfrm}[w]$ তখনই এবং কেবল তখনই, যখন
    $R_a ww'$-যুক্ত প্রত্যেক $w' \in W$-তে $\mSat{M}{!B}[w']$।}
  \end{enumerate}
\end{defn}

এখানেই আমাদের অভিগম্যতা-সম্পর্কের ওপর বিধিনিষেধ নিয়ে ভাবতে হয়। কারণ \olref{defn:sub:mmodels-box} অনুসারে $R_a ww'$-যুক্ত কোনো~$w'$ না থাকলেও !!a{formula}~$\Knows_a !B$ জগৎ~$w$-তে সত্য। স্বাভাবিক মোডাল যুক্তিবিদ্যাতেও একই শর্ত: কোনো জগতের পরবর্তী অভিগম্য জগৎ না থাকলে সব $\Box$-সূত্র সেখানে শূন্যতাবশে সত্য। কিন্তু $\Knows$-কে \emph{জ্ঞান} হিসেবে ধরলে এটি অত্যন্ত অস্বাভাবিক মনে হয়। বস্তুত, আমরা সাধারণত মনে করি না যে কোনো পরিস্থিতিতে একজন কর্তা একই সঙ্গে $!A$ এবং $\lnot !A$—উভয়ই জানতে পারে।

এর একটি সমাধান হলো জ্ঞানতাত্ত্বিক যুক্তিবিদ্যায় অভিগম্যতা-সম্পর্ককে সব সময় \emph{প্রতিবিম্বী} রাখা। মোটামুটিভাবে এর অর্থ, প্রকৃত জগৎ কর্তার তথ্যের সঙ্গে সঙ্গতিপূর্ণ। বাস্তবে জ্ঞানতাত্ত্বিক যুক্তিবিদ্যায় প্রায়ই S5 ব্যবহার করা হয়; তবে $\Knows_a$ সম্পর্ক দিয়ে ঠিক কী বোঝাতে চাই, তার ওপর নির্ভর করে দুর্বলতর পদ্ধতিও ব্যবহৃত হতে পারে।

\begin{figure}
  \begin{center}
    \begin{tikzpicture}[modal]
      \node[world] (w1) [label={[align=right]right:\mTrue{p}\\ \mFalse{q}}]{$w_1$};
      \node[world] (w2) [label={[align=right]right:\mTrue{p}\\ \mTrue{q}},
        above right=of w1]{$w_2$};
      \node[world] (w3) [label={[align=right]right:\mFalse{p}\\ \mFalse{q}},
        below right=of w1] {$w_3$};
      \draw[<->] (w1) to node {$a$} (w2);
      \draw[->,loop left] (w1) to node {$a,b$} (w1);
      \draw[<->] (w1) to [swap] node {$b$} (w3);
      \draw[->,loop above] (w2) to node {$a,b$} (w2) ;
      \draw[->,loop below] (w3) to node {$a,b$} (w3) ;
    \end{tikzpicture}
  \end{center}
  \caption{একটি সরল জ্ঞানতাত্ত্বিক মডেল।}
  \ollabel{fig:simple}
\end{figure}

\begin{prob}
  \olref[aml][el][trw]{fig:simple}-এ নিচের কোনগুলি সত্য, বিবেচনা করুন:
  \begin{enumerate}
    \item $\mSat{M}{\lnot q}[w_1]$;
    \item $\mSat{M}{\Knows_a \lnot q}[w_1]$;
    \item $\mSat{M}{\Knows_b \lnot q}[w_1]$;
    \item $\mSat{M}{\Knows_b q \lor \Knows_b \lnot q}[w_2]$;
    \item $\mSat{M}{\Knows_a( \Knows_b q \lor \Knows_b \lnot q)}[w_2]$;
    \item $\mSat{M}{\EKnows_{\{a,b\}} \lnot q}[w_3]$;
    \end{enumerate}
\end{prob}

কোনো জগতে সত্যতার মৌলিক সংজ্ঞা দেওয়া হলো। এখন স্বাভাবিক মোডাল যুক্তিবিদ্যার মোডাল বৈধতা ও অনুসিদ্ধান্তের মতো অন্য অর্থতাত্ত্বিক ধারণাগুলিও মোডাল অপারেটরগুলির এই নতুন ব্যাখ্যায় সরাসরি প্রযোজ্য।

এবার সর্বস্তরে পারস্পরিক জ্ঞানের অপারেটর~$\CKnows_G$-এর সত্যতার শর্তও দিতে পারি। \olref[sfr][rel][ops]{sec} থেকে স্মরণ করুন, সম্পর্ক~$R$-এর \emph{পরিযায়ী আবরণ}~$R^+$ সংজ্ঞায়িত হয় এভাবে:
\begin{align*}
  R^+ &= \bigcup_{ n \in \mathbb{N}} R^n,
\intertext{যেখানে}
  R^0 & = R \text{ এবং}\\
  R^{n+1} & = \Setabs{\tuple{x, z}}{\exists y (R^n xy \land Ryz)}.
\end{align*}
তাহলে $G$ যদি কর্তাদের একটি গোষ্ঠী হয়, সব কর্তার অভিগম্যতা-সম্পর্কের সংযুক্তির পরিযায়ী আবরণ হিসেবে $R_G = ( \bigcup_{b
\in G} R_b )^+$ সংজ্ঞায়িত করি।

\begin{defn}
$G' \subseteq G$ হলে $\mSat{M}{\CKnows_{G'} !A}[ w]$ তখনই এবং কেবল তখনই, যখন $R_{G'} w w'$ হয় এমন প্রত্যেক $w'$-তে $\mSat{M}{!A}[w']$।
\end{defn}

\end{document}
